%% kapa: exam=1 round=34 year=2023 period=2 form=A subject=appraisal_law subject_ko=감정평가관계법규 seq=38 number=38
%% source: Q-Net 감정평가사 기출문제 (cst00309) artlSeq=5211789 file=raw/1008/5211789/2023년도 제34회 감정평가사 1차시험 2교시.pdf
%% conversion: kordoc 4.24.51 -> tools/gen1.py (auto); review=proofread-sonnet; flags=none
\begin{question}{38}
도시 및 주거환경정비법령상 정비사업에 관한 설명이다. (　　)에 들어갈 내용으로 옳은 것은?
\begin{framedtext}
○ 시장·군수등이 아닌 사업시행자가 정비사업을 시행하려는 경우에는 사업시행계획서에 정관등과 그 밖에 국토교통부령으로 정하는 서류를 첨부하여 시장·군수등에게 제출하고 사업시행계획(ㄱ)을(를) 받아야 한다. \par
○ 시장·군수등이 아닌 사업시행자가 정비사업 공사를 완료한 때에는 대통령령으로 정하는 방법 및 절차에 따라 시장·군수등의 준공(ㄴ)을(를) 받아야 한다. \par
\end{framedtext}
\begin{choices}
  \item ㄱ: 승인, ㄴ: 인가
  \item ㄱ: 승인, ㄴ: 허가
  \item ㄱ: 인가, ㄴ: 승인
  \item ㄱ: 인가, ㄴ: 인가
  \item ㄱ: 인가, ㄴ: 허가
\end{choices}
\end{question}
\answer{4}
